\section{Three rows}\label{sec:three}

Take $m\ge3$ and number the rows $1,2,3$ downward. We write a column as a
letter. The singletons are $a=\{1\}$, $b=\{2\}$, and $c=\{3\}$, the doubles are
$A=\{1,2\}$, $B=\{1,3\}$, and $C=\{2,3\}$, and the triple is $X=\{1,2,3\}$. We
call $a$ and $c$ the \emph{outer} letters. A blackout without empty columns is
then a word in these $7$ letters. We also write $A$, $B$, $C$ for the
corresponding pairs of rows, so that the pair set $F$ is a set of letters.

\subsection{The candidates and the window rules}

\begin{proposition}[the two types]\label{prop:threeclass}
Let $m\ge3$.
\begin{enumerate}[(a)]
\item A valid blackout of the $3\times m$ grid has at most $m+2$ points. If it has
$m+2$ points, then its word consists of singletons and either one $X$
\textup{(type~I)} or two different doubles \textup{(type~II)}.
\item For every word of type~I or~II, the graph $G_S$ is a tree. Hence the
blackout is axis-valid and hides no axis-aligned rectangle.
\item There are $m3^{m-1}+\binom m2\cdot6\cdot3^{m-2}=m^23^{m-1}$ words of type~I
or~II.
\end{enumerate}
\end{proposition}

\begin{proof}
(a) Let $S$ be valid. By Lemma~\ref{lem:girth}, $G_S$ has no $4$-cycle and no
$6$-cycle. So no $2$ columns share $2$ black rows. Suppose first that some
column is a triple. Then every other column has at most $1$ black point, and
$|S|\le3+(m-1)$, with equality only for type~I. Suppose next that no column is a
triple. The doubles have different row pairs. Three doubles would carry
all $3$ pairs, and their columns would form a $6$-cycle with the $3$ rows. So
there are at most $2$ doubles and $|S|\le4+(m-2)$, with equality only for
type~II.

(b) The nonsingleton columns and the $3$ rows form a path or a star, and each
singleton column is a leaf. Apply Lemma~\ref{lem:girth}. A hidden axis-aligned
rectangle would be a $4$-cycle.

(c) Place the triple, or an ordered pair of different doubles in $2$ of the
columns, and choose each singleton freely.
\end{proof}

We call the words of type~I or~II \emph{candidates}. The only tilted
rectangles are the diamonds $D(i,1)$ (Lemma~\ref{lem:tilted-narrow}). We say
that the diamond $D(i,1)$ \emph{fails} in a candidate if it shows $1$ or $2$
corners and this presentation has a rival in the sense of
Lemma~\ref{lem:rival}. Otherwise it \emph{passes}.

\begin{lemma}[diamond test]\label{lem:diamond}
A candidate is valid if and only if every diamond $D(i,1)$ passes.
\end{lemma}

\begin{proof}
Axis-aligned rectangles have distinct presentations by
Proposition~\ref{prop:threeclass}(b). Two diamonds $D(i,1)$ and $D(i',1)$ with
$i<i'$ share at most $1$ corner, which is a point of row $2$. If they had the
same presentation, the points of rows $1$ and $3$ in the columns $i+1$ and
$i'+1$ would be black by Lemma~\ref{lem:diff}. Then both columns would be $B$
or $X$, and a candidate has at most $1$ such column. So the candidate is valid
if and only if no diamond has the presentation of an axis-aligned rectangle.
A diamond that shows at least $3$ corners is safe by
Lemma~\ref{lem:threecorners}. A hidden diamond is safe because no axis-aligned
rectangle is hidden. The remaining diamonds are the subject of
Lemma~\ref{lem:rival}.
\end{proof}

A \emph{window} is a factor of $3$ consecutive letters $\ell\,t\,r$ of a
candidate, in the columns $i$, $i+1$, and $i+2$. It \emph{passes} or
\emph{fails} with $D(i,1)$. The corners of $D(i,1)$ are the \emph{tips} $(i,2)$
and $(i+2,2)$ and the points $(i+1,1)$ and $(i+1,3)$ of the middle column. A
tip is black exactly when its letter is $b$, $A$, $C$, or $X$.

\begin{lemma}[window rules]\label{lem:windows}
Let $\ell\,t\,r$ be a window of a candidate with pair set $F$.
\begin{enumerate}[(a)]
\item If $t$ is a singleton, the window fails if and only if both tips are black and
$B\in F$.
\item If $t=A$, the window fails if and only if both tips are black, or exactly $1$
tip is black and the letter of the white tip is $c$ or $B$. If $t=C$, the same
holds with $a$ in place of $c$.
\item If $t$ is $B$ or $X$, the window fails if and only if both tips are white
with the same letter, or exactly $1$ tip is black and the letter of the white tip
is $a$ with $A\in F$ or $c$ with $C\in F$.
\end{enumerate}
\end{lemma}

\begin{proof}
(a) At most $1$ point of the middle column is black. If both tips are white, at
least $3$ corners show. Suppose that exactly $1$ tip is black and the other tip
$(x,2)$ is white. For $t=b$ again $3$ corners show. For $t=a$ the diamond shows
the oblique pair $(x,2)$, $(i+1,3)$. A rival would need $(i+1,2)$ to be black
(Lemma~\ref{lem:rival}(d)), and it is white. The case $t=c$ is the same.
Suppose finally that both tips are black. For $t=b$ the diamond shows the
vertical pair in rows $1$ and $3$, which has a rival if and only if $B\in F$
(Lemma~\ref{lem:rival}(b)). For $t=a$ it shows only $(i+1,3)$, and
$B_{i+1}=\{1\}$. By Lemma~\ref{lem:rival}(a) there is a rival if and only if
$B\in F$. The case $t=c$ is the same.

(b) Let $t=A$. Then $(i+1,1)$ and $(i+1,2)$ are black and $(i+1,3)$ is white. If
both tips are white, $3$ corners show. If both are black, the diamond shows only
$(i+1,3)$. The candidate has type~II, so $F$ contains a second pair, and this
pair contains row $3$. Its other row is black in column $i+1$, so there is a
rival by Lemma~\ref{lem:rival}(a). If exactly $1$ tip is black and $(x,2)$ is
the white tip, the diamond shows the oblique pair $(x,2)$, $(i+1,3)$. Since
$(i+1,2)$ is black, there is a rival if and only if $(x,3)$ is black, which
means that the letter of column $x$ is $c$ or $B$. The case $t=C$ follows by
reflecting the rows.

(c) Now both points of the middle column are black. If both tips are black, the
diamond is hidden and passes. A white tip has the letter $a$ or $c$, because the
letters $A$, $C$ have black tips and the word has only one letter $B$ or $X$.
If both tips are white, the diamond shows a horizontal pair, and by
Lemma~\ref{lem:rival}(c) there is a rival if and only if the $2$ letters are
equal. If exactly $1$ tip is black, the diamond shows only the white tip
$(x,2)$, and $B_x=\{s\}$ with $s=1$ or $s=3$. By Lemma~\ref{lem:rival}(a) there
is a rival if and only if $\{2,s\}\in F$.
\end{proof}

\begin{corollary}\label{cor:threevalue}
For $m\ge3$ the word $Xa^{m-1}$ is valid, so $f(3,m)=m+2$, and the optima of
the $3\times m$ grid are the candidates in which every window passes.
\end{corollary}

\begin{proof}
In $Xa^{m-1}$ every window has the middle letter $a$ and a white right tip.
It passes by Lemma~\ref{lem:windows}(a). Now use
Proposition~\ref{prop:threeclass}(a) and Lemma~\ref{lem:diamond}.
\end{proof}

We sort the candidates into $7$ \emph{classes} by their nonsingleton letters
read from left to right, namely $X$, $AC$, $CA$, $AB$, $BA$, $BC$, and $CB$. We
write $AC_m$ for the set of valid words of width $m$ in the class $AC$, and
likewise for the other classes. The row reflection exchanges $a$ with $c$ and
$A$ with $C$, and it fixes $b$, $B$, and $X$. The column reversal reverses the
word. Both preserve validity. Hence
\begin{equation}\label{eq:classsym}
|AC_m|=|CA_m|\qquad\text{and}\qquad|AB_m|=|BA_m|=|CB_m|=|BC_m| .
\end{equation}
Table~\ref{tab:windows} lists the failing windows in the classes $AC$, $AB$, and
$X$. It is read off from Lemma~\ref{lem:windows}.

\begin{figure}[t]
\begin{minipage}[b]{0.5\textwidth}
\centering\footnotesize
\begin{tabular}{@{}lll@{}}
\toprule
class & middle & failing windows \\
\midrule
$AC$ & $A$ & $bAb$, $bAc$, $cAb$, $bAC$, $cAC$ \\
 & $C$ & $bCb$, $bCa$, $aCb$, $ACb$, $ACa$ \\
\midrule
$AB$ & singleton $s$ & $bsb$, $bsA$, $Asb$ \\
 & $A$ & $bAb$, $bAc$, $cAb$, $bAB$ \\
 & $B$ & $aBa$, $cBc$, $aBb$, $bBa$, $ABa$ \\
\midrule
$X$ & singleton $s$ & $bsb$, $bsX$, $Xsb$ \\
 & $X$ & $aXa$, $cXc$, $aXb$, $bXa$, $bXc$, $cXb$ \\
\bottomrule
\end{tabular}
\captionof{table}{The failing windows, by class and middle letter. Here $s$ is any
singleton. In the class $AC$ every window with a singleton in the middle passes.}
\label{tab:windows}
\end{minipage}\hfill
\begin{minipage}[b]{0.46\textwidth}
\centering
\resizebox{\textwidth}{!}{\begin{tikzpicture}[x=0.62cm,y=0.62cm,baseline=(current bounding box.north)]
  \begin{scope}[xshift=0.000cm]
    \draw[black,line width=0.25pt,dash pattern=on 0.93pt off 0.93pt] (1,-1) -- (5,-1) (1,-2) -- (5,-2) (1,-3) -- (5,-3) (1,-1) -- (1,-3) (2,-1) -- (2,-3) (3,-1) -- (3,-3) (4,-1) -- (4,-3) (5,-1) -- (5,-3);
    \draw[kred,line width=1.76pt,line join=round] (1,-2) -- (2,-1) -- (3,-2) -- (2,-3) -- cycle;
    \filldraw[fill=kgray,draw=black,line width=0.37pt] (1,-1) circle (0.3);
    \draw[black,line width=0.37pt] (0.788,-0.788) -- (1.212,-1.212) (0.788,-1.212) -- (1.212,-0.788);
    \filldraw[fill=kgray,draw=black,line width=0.37pt] (2,-1) circle (0.3);
    \draw[black,line width=0.37pt] (1.788,-0.788) -- (2.212,-1.212) (1.788,-1.212) -- (2.212,-0.788);
    \filldraw[fill=white,draw=black,line width=0.37pt] (3,-1) circle (0.3);
    \filldraw[fill=white,draw=black,line width=0.37pt] (4,-1) circle (0.3);
    \filldraw[fill=white,draw=black,line width=0.37pt] (5,-1) circle (0.3);
    \filldraw[fill=kyellow,draw=black,line width=0.37pt] (1,-2) circle (0.3);
    \filldraw[fill=kgray,draw=black,line width=0.37pt] (2,-2) circle (0.3);
    \draw[black,line width=0.37pt] (1.788,-1.788) -- (2.212,-2.212) (1.788,-2.212) -- (2.212,-1.788);
    \filldraw[fill=kgray,draw=black,line width=0.37pt] (3,-2) circle (0.3);
    \draw[black,line width=0.37pt] (2.788,-1.788) -- (3.212,-2.212) (2.788,-2.212) -- (3.212,-1.788);
    \filldraw[fill=kgray,draw=black,line width=0.37pt] (4,-2) circle (0.3);
    \draw[black,line width=0.37pt] (3.788,-1.788) -- (4.212,-2.212) (3.788,-2.212) -- (4.212,-1.788);
    \filldraw[fill=white,draw=black,line width=0.37pt] (5,-2) circle (0.3);
    \filldraw[fill=white,draw=black,line width=0.37pt] (1,-3) circle (0.3);
    \filldraw[fill=kyellow,draw=black,line width=0.37pt] (2,-3) circle (0.3);
    \filldraw[fill=white,draw=black,line width=0.37pt] (3,-3) circle (0.3);
    \filldraw[fill=kgray,draw=black,line width=0.37pt] (4,-3) circle (0.3);
    \draw[black,line width=0.37pt] (3.788,-2.788) -- (4.212,-3.212) (3.788,-3.212) -- (4.212,-2.788);
    \filldraw[fill=kgray,draw=black,line width=0.37pt] (5,-3) circle (0.3);
    \draw[black,line width=0.37pt] (4.788,-2.788) -- (5.212,-3.212) (4.788,-3.212) -- (5.212,-2.788);
    \node[lab] at (1,-0.45) {1};
    \node[lab] at (2,-0.45) {2};
    \node[lab] at (3,-0.45) {3};
    \node[lab] at (4,-0.45) {4};
    \node[lab] at (5,-0.45) {5};
    \node[lab] at (0.45,-1) {1};
    \node[lab] at (0.45,-2) {2};
    \node[lab] at (0.45,-3) {3};
    \node[capt,anchor=north] at (3.00,-3.45) {(a) $aAbCc$ is valid};
  \end{scope}
  \begin{scope}[xshift=4.092cm]
    \draw[black,line width=0.25pt,dash pattern=on 0.93pt off 0.93pt] (1,-1) -- (5,-1) (1,-2) -- (5,-2) (1,-3) -- (5,-3) (1,-1) -- (1,-3) (2,-1) -- (2,-3) (3,-1) -- (3,-3) (4,-1) -- (4,-3) (5,-1) -- (5,-3);
    \draw[kblue,line width=1.76pt,line join=round] (2,-2) -- (4,-2) -- (4,-3) -- (2,-3) -- cycle;
    \draw[kred,line width=1.76pt,line join=round] (1,-2) -- (2,-1) -- (3,-2) -- (2,-3) -- cycle;
    \filldraw[fill=white,draw=black,line width=0.37pt] (1,-1) circle (0.3);
    \filldraw[fill=kgray,draw=black,line width=0.37pt] (2,-1) circle (0.3);
    \draw[black,line width=0.37pt] (1.788,-0.788) -- (2.212,-1.212) (1.788,-1.212) -- (2.212,-0.788);
    \filldraw[fill=white,draw=black,line width=0.37pt] (3,-1) circle (0.3);
    \filldraw[fill=white,draw=black,line width=0.37pt] (4,-1) circle (0.3);
    \filldraw[fill=white,draw=black,line width=0.37pt] (5,-1) circle (0.3);
    \filldraw[fill=kgray,draw=black,line width=0.37pt] (1,-2) circle (0.3);
    \draw[black,line width=0.37pt] (0.788,-1.788) -- (1.212,-2.212) (0.788,-2.212) -- (1.212,-1.788);
    \filldraw[fill=kgray,draw=black,line width=0.37pt] (2,-2) circle (0.3);
    \draw[black,line width=0.37pt] (1.788,-1.788) -- (2.212,-2.212) (1.788,-2.212) -- (2.212,-1.788);
    \filldraw[fill=kgray,draw=black,line width=0.37pt] (3,-2) circle (0.3);
    \draw[black,line width=0.37pt] (2.788,-1.788) -- (3.212,-2.212) (2.788,-2.212) -- (3.212,-1.788);
    \filldraw[fill=kgray,draw=black,line width=0.37pt] (4,-2) circle (0.3);
    \draw[black,line width=0.37pt] (3.788,-1.788) -- (4.212,-2.212) (3.788,-2.212) -- (4.212,-1.788);
    \filldraw[fill=white,draw=black,line width=0.37pt] (5,-2) circle (0.3);
    \filldraw[fill=white,draw=black,line width=0.37pt] (1,-3) circle (0.3);
    \filldraw[fill=kyellow,draw=black,line width=0.37pt] (2,-3) circle (0.3);
    \filldraw[fill=white,draw=black,line width=0.37pt] (3,-3) circle (0.3);
    \filldraw[fill=kgray,draw=black,line width=0.37pt] (4,-3) circle (0.3);
    \draw[black,line width=0.37pt] (3.788,-2.788) -- (4.212,-3.212) (3.788,-3.212) -- (4.212,-2.788);
    \filldraw[fill=kgray,draw=black,line width=0.37pt] (5,-3) circle (0.3);
    \draw[black,line width=0.37pt] (4.788,-2.788) -- (5.212,-3.212) (4.788,-3.212) -- (5.212,-2.788);
    \node[lab] at (1,-0.45) {1};
    \node[lab] at (2,-0.45) {2};
    \node[lab] at (3,-0.45) {3};
    \node[lab] at (4,-0.45) {4};
    \node[lab] at (5,-0.45) {5};
    \node[lab] at (0.45,-1) {1};
    \node[lab] at (0.45,-2) {2};
    \node[lab] at (0.45,-3) {3};
    \node[capt,anchor=north] at (3.00,-3.45) {(b) $bAbCc$ fails};
  \end{scope}
\end{tikzpicture}
}
\caption{The diamond test. (a) In $aAbCc$ the diamond $D(1,1)$ shows $(1,2)$
and $(2,3)$. This pair has no rival, since $(1,3)$ is white. (b) In $bAbCc$ it
shows only $(2,3)$, and so does the blue rectangle. The window $bAb$ fails.}
\label{fig:diamond}
\end{minipage}
\end{figure}

\subsection{\texorpdfstring{The class $AC$}{The class AC}}

Set $G_{AC}(z)=\sum_{m\ge2}|AC_m|z^m$, and define $G_{AB}$ and $G_X$ in the same
way, where the sum for $G_X$ starts at $m=1$. Write $\Phi=1-3z$. The nonempty
singleton strings with a given first letter have the series $z/\Phi$, and so do
those with a given last letter.

\begin{proposition}[the class $AC$]\label{prop:threeGAC}
We have
\begin{equation}\label{eq:GAC}
G_{AC}(z)=\frac{z^2-4z^3+10z^4-8z^5-3z^6}{(1-3z)^3},
\end{equation}
and
\begin{equation}\label{eq:GACcoef}
|AC_m|=2(m-1)(3m+2)\,3^{m-5}\qquad(m\ge4).
\end{equation}
\end{proposition}

\begin{proof}
Write a word of the class as $UAVCW$ with singleton strings $U$, $V$, and $W$.
By Table~\ref{tab:windows} only the windows centered on $A$ and on $C$ matter.
Read $C$ as $b$ when it is a tip, and likewise $A$. Then the window $sAt$ fails
exactly for $st\in\{bb,bc,cb\}$, and $sCt$ fails exactly for
$st\in\{bb,ba,ab\}$. We count by the length of $V$.

If $V$ is empty, a nonempty $U$ must end in $a$ and a nonempty $W$ must start
with $c$. This case contributes $z^2(1+z/\Phi)^2$.

If $V$ is a single letter $t$, let $p_t$ be the number of letters allowed at the
end of $U$ and $q_t$ the number allowed at the start of $W$. Then
$(p_a,p_b,p_c)=(3,1,2)$ and $(q_a,q_b,q_c)=(2,1,3)$. The contribution is
\[
z^3\sum_t\Bigl(1+\frac{p_tz}{\Phi}\Bigr)\Bigl(1+\frac{q_tz}{\Phi}\Bigr)
=z^3\Bigl(3+\frac{12z}{\Phi}+\frac{13z^2}{\Phi^2}\Bigr).
\]

If $V$ has at least $2$ letters, its first letter and the last letter of $U$
form one of $6$ allowed pairs, or $U$ is empty and the first letter of $V$ is
free. This gives the factor $3+6z/\Phi$. The other end gives the same factor,
and the letters of $V$ between its first and last are free. The contribution
is $z^4(3+6z/\Phi)^2/\Phi=9z^4(1-z)^2/\Phi^3$.

Adding the $3$ cases gives~\eqref{eq:GAC}. Since
$[z^m]z^k\Phi^{-3}=\binom{m-k+2}23^{m-k}$, the coefficient of $z^m$ for $m\ge4$
is~\eqref{eq:GACcoef}.
\end{proof}

Let $E_m$ be the number of optima of the $3\times m$ grid with a column that is
black in rows $1$ and $3$ (\seqnum{A398448}). Such a column is $B$ or $X$. So
$E_m$ counts the optima in the $5$ classes other than $AC$ and $CA$, and
by~\eqref{eq:classsym},
\begin{equation}\label{eq:classsum}
a_3(m)=2|AC_m|+4|AB_m|+|X_m|=2|AC_m|+E_m\qquad(m\ge3).
\end{equation}

\subsection{\texorpdfstring{The classes $AB$ and $X$}{The classes AB and X}}\label{sec:threeruns}

Put $L=1-2z-2z^2$, $N=1+2z^2$, and $\Psi=LN=1-2z-4z^3-4z^4$. Call a singleton
string \emph{restricted} if it has no $2$ letters $b$ at distance $2$. In the
classes $AB$ and $X$ the windows $bsb$ fail, so every run of singletons is
restricted. Give the letter $b$ the type $1$ and the outer letters the type
$0$. For types $i$ and $j$, let $h_{ij}(z)$ count by length the strings
$y_1\cdots y_n$ with $n\ge0$ that can follow $2$ letters of types $i$ and $j$ in
a restricted string. Write $h_{ij}=h_{ij}^0+h_{ij}^1$, where $h_{ij}^\beta$
counts the strings whose last letter has type $\beta$, and the empty string is
given the type $j$. The letter $y_1$ can be $a$ or $c$, and it can be $b$ when
$i=0$. Hence, for $\beta\in\{0,1\}$,
\[
h_{ij}^\beta=[j=\beta]+2z\,h_{j0}^\beta+[i=0]\,z\,h_{j1}^\beta .
\]
Solving these $2$ linear systems gives
\begin{equation}\label{eq:hsol}
\begin{aligned}
\Psi h_{00}^0&=1+2z^2+2z^3,&\Psi h_{00}^1&=z+z^2,&\Psi h_{10}^0&=1,&\Psi h_{10}^1&=2z^2+2z^3,\\
\Psi h_{00}&=1+z+3z^2+2z^3,&\Psi h_{01}&=1+z,&\Psi h_{10}&=1+2z^2+2z^3,&\Psi h_{11}&=1.
\end{aligned}
\end{equation}

We read a run of singletons away from the double or triple next to it, as
$w_1w_2\cdots w_n$ with $n\ge1$. We need $4$ series for such runs. Let $g_0$
count the restricted runs whose first letter is a given outer letter, and
$g_1$ those with $w_1=b$. Let $u$ and $v$ count the same runs under the extra
condition that $w_2$, if present, is outer. A missing letter before $w_1$ acts
as a letter of type $0$, and the extra condition acts as a letter of type
$1$. Hence
\begin{equation}\label{eq:runs}
g_0=zh_{00},\qquad g_1=zh_{01},\qquad u=zh_{10},\qquad v=zh_{11}.
\end{equation}

\begin{proposition}[the class $X$]\label{prop:GX}
We have
\[
G_X(z)=z\bigl(1+4u+2v+2u^2+v^2\bigr)
=\frac{z(1+2z-7z^2-4z^3-2z^4-8z^5-4z^6)}{L^2N}.
\]
\end{proposition}

\begin{proof}
Write the word as $UXW$. Both runs are restricted. The failing windows $bsX$
and $Xsb$ of Table~\ref{tab:windows} say that the second letter from $X$ is
outer. The window centered on $X$ allows the neighbor pairs $ac$, $ca$, and
$bb$. So a run on one side only contributes $2u+v$, for either side, and
runs on both sides contribute $2u^2+v^2$. The second form follows
from~\eqref{eq:hsol} and~\eqref{eq:runs}.
\end{proof}

Now write a word of the class $AB$ as $UAVBW$. For a letter $t$ next to $A$ in
$V$, let $\ell_t$ count the strings $U$ that are allowed, including the empty
one. For a letter $t$ next to $B$ in $V$, let $r_t$ count the strings $W$ that
are allowed. By Table~\ref{tab:windows}, the window $bsA$ makes the second
letter of $U$ from $A$ outer, and nothing restricts the second letter of
$W$. The windows centered on $A$ allow the letters $a,b,c$ before $At$ for
$t=a$, only $a$ for $t=b$, and $a,c$ for $t=c$. The windows centered on $B$
allow only $c$ after $tB$ for $t=a$, the letters $b,c$ for $t=b$, and $a,b$
for $t=c$. Hence
\begin{equation}\label{eq:lr}
(\ell_a,\ell_b,\ell_c)=(1+2u+v,\ 1+u,\ 1+2u),\qquad
(r_a,r_b,r_c)=(1+g_0,\ 1+g_0+g_1,\ 1+g_0+g_1).
\end{equation}

\begin{proposition}[the class $AB$, $P_B$ by computer algebra]\label{prop:GAB}
With the series of~\eqref{eq:hsol} to~\eqref{eq:lr},
\begin{equation}\label{eq:GAB}
\begin{aligned}
G_{AB}(z)={}&z^2(1+2u)(1+g_0+g_1)+z^3(\ell_ar_a+\ell_br_b+\ell_cr_c)\\
&+z^4\sum_{\sigma\in\{0,1\}}\lambda_\sigma\bigl(h^0_{\sigma0}\,(r_a+r_c)+2h^1_{\sigma0}\,r_b\bigr),
\end{aligned}
\end{equation}
where $\lambda_0=\ell_a+\ell_c$ and $\lambda_1=\ell_b$. Moreover $G_{AB}=P_B/\Psi^3$
with
\begin{equation}\label{eq:PB}
P_B(z)=z^2(1+z)(1+z^2-3z^3+4z^6+4z^7).
\end{equation}
\end{proposition}

\begin{proof}
Let $k$ be the length of $V=v_1\cdots v_k$. For $k=0$ the windows $sAB$ and
$ABt$ of Table~\ref{tab:windows} require $s\ne b$ and $t\ne a$, which gives the
first term. For $k=1$ the window $Av_1B$ passes, and the windows centered on
$A$ and $B$ give $\ell_tr_t$ for $v_1=t$.

Let $k\ge2$. The window $Av_1v_2$ makes $v_2$ outer, and the window
$v_{k-1}v_kB$ passes. So $V$ is a restricted string whose second letter is
outer, and its weight is $\ell_{v_1}r_{v_k}$. Let $\sigma$ be the type of
$v_1$. Summing $\ell_{v_1}$ over the letters $v_1$ of type $\sigma$ gives
$\lambda_\sigma$. The string $v_3\cdots v_k$ follows $2$ letters of types
$\sigma$ and $0$. If it is empty, then $v_k=v_2$ is outer and the $2$ choices
give $r_a+r_c$. If it is not empty and $v_k$ is outer, there are $2$ choices
for $v_2$, and $h^0_{\sigma0}$ counts each string $v_3\cdots v_{k-1}$ twice,
once for each $v_k$. So the weight is again $h^0_{\sigma0}(r_a+r_c)$. If
$v_k=b$, the $2$ choices for $v_2$ give $2h^1_{\sigma0}r_b$. This
proves~\eqref{eq:GAB}. Substituting~\eqref{eq:hsol} and~\eqref{eq:runs} and
expanding gives~\eqref{eq:PB}.
\end{proof}

\subsection{Generating function, closed form, and asymptotics}

\begin{theorem}[three-row generating function, $P$ by computer algebra]\label{thm:threegf}
\leavevmode
\begin{enumerate}[(a)]
\item We have
\begin{equation}\label{eq:PQ}
\sum_{m\ge1}a_3(m)z^m=2G_{AC}+4G_{AB}+G_X=\frac{P(z)}{Q(z)},
\end{equation}
where $Q=(1-3z)^3(1-2z-2z^2)^3(1+2z^2)^3$ and
\begin{align*}
P(z)={}&z-3z^2-34z^3+228z^4-641z^5+1307z^6-2186z^7\\
&+2056z^8-236z^9-1060z^{10}+2584z^{11}-4264z^{12}\\
&-1440z^{13}+704z^{14}-1024z^{15}+2080z^{16}+2176z^{17}+384z^{18}.
\end{align*}
\item The polynomials $P$ and $Q$ are coprime. So $Q$ is the reduced denominator,
and $a_3(m)$ satisfies the linear recurrence of order $15$ with characteristic
polynomial $x^{15}Q(1/x)$ for $m\ge19$ and no recurrence of smaller order.
\item We have
\[
\sum_{m\ge3}E_mz^m=\frac{z^3\,P_E(z)}{(1-2z-2z^2)^3(1+2z^2)^3},
\]
where $P_E=43-90z+128z^2-344z^3+44z^4+80z^5+328z^6+1728z^7+1888z^8+2208z^9
+2016z^{10}+640z^{11}$.
\end{enumerate}
\end{theorem}

\begin{proof}
(a) For $m\ge3$ the coefficient of $z^m$ in $2G_{AC}+4G_{AB}+G_X$ is $a_3(m)$
by~\eqref{eq:classsum}. The coefficients of $z$ and $z^2$ are $1$ and $12$.
These are $a_3(1)$ and $a_3(2)=a_2(3)$, by Section~\ref{sec:setup} and
Theorem~\ref{thm:strip}. Bringing the $3$ series of
Propositions~\ref{prop:threeGAC}, \ref{prop:GX}, and~\ref{prop:GAB} to the common
denominator $Q$ gives $P$.

(b) The factors $\Phi$, $L$, and $N$ of $Q$ are irreducible and distinct. At
$z=\frac13$ only $G_{AC}$ has a pole, and the numerator of~\eqref{eq:GAC} takes
the value $\frac4{81}$ there. So the pole at $\frac13$ has order $3$. The remainders of
$P_B$ modulo $L$ and modulo $N$ are $\frac z2$ and $\frac14-\frac z2$. They are
not zero, so $G_{AB}$ has poles of order $3$ at the roots of $L$ and of $N$.
The series $G_X$ has poles of order at most $2$ there, and $G_{AC}$ has none.
Hence the series~\eqref{eq:PQ} has a pole of order $3$ at every root of $Q$, and
$P$ and $Q$ are coprime. The statement on recurrences follows, since
$\deg P=18$.

(c) By~\eqref{eq:classsum} the series is $4G_{AB}+G_X$ without its terms in
$z$ and $z^2$, which are $z+10z^2$.
\end{proof}

The polynomials $P$, $P_B$, and $P_E$ were expanded in exact arithmetic with
SymPy~\cite{sympy2017}.

Let $\rho=1+\sqrt3$ and $\sigma=1-\sqrt3$, so that $L=(1-\rho z)(1-\sigma z)$.

\begin{theorem}[the number of three-row optima, checked by computer algebra]\label{thm:threecount}
For every $m\ge4$,
\begin{equation}\label{eq:threecount}
a_3(m)=4(m-1)(3m+2)\,3^{m-5}+p_+(m)\,\rho^m+p_-(m)\,\sigma^m+(-2)^{\lfloor m/2\rfloor}\kappa(m),
\end{equation}
where
\[
432\,p_\pm(m)=(19\pm11\sqrt3)m^2+(61\pm41\sqrt3)m+104\pm28\sqrt3,
\]
and $\kappa(m)=(2m^2+41m-104)/216$ for even $m$ and
$\kappa(m)=(7m^2-32m-175)/432$ for odd $m$. The last $3$ terms are $E_m$, and
\begin{equation}\label{eq:Eclosed}
E_m=p_+(m)\,\rho^m+p_-(m)\,\sigma^m+(-2)^{\lfloor m/2\rfloor}\kappa(m)
\end{equation}
holds for every $m\ge2$. Moreover $E_m>0$,
\begin{equation}\label{eq:asymp1}
E_m\sim\frac{19+11\sqrt3}{432}\,m^2(1+\sqrt3)^m,\qquad
a_3(m)=\frac4{81}m^23^m-\frac4{243}(m+2)3^m+O\bigl(m^2(1+\sqrt3)^m\bigr),
\end{equation}
and the fraction of the $m^23^{m-1}$ candidates that are valid tends to
$\frac4{27}$.
\end{theorem}

\begin{proof}
By~\eqref{eq:classsum} and~\eqref{eq:GACcoef} we have
$a_3(m)=4(m-1)(3m+2)3^{m-5}+E_m$ for $m\ge4$. So it is enough to
prove~\eqref{eq:Eclosed}. By~\eqref{eq:classsum}, $E_m$ is the coefficient of
$z^m$ in $4G_{AB}+G_X$ for $m\ge3$. This also holds for $m=2$. Indeed the optima
of the $3\times2$ grid are the transposes of the optima of
Theorem~\ref{thm:strip}. They are the $12$ words of type~I or~II, and $10$ of
them contain $B$ or $X$. By Propositions~\ref{prop:GX} and~\ref{prop:GAB}, the
series $4G_{AB}+G_X$ has the denominator $L^3N^3$ of degree $12$ and a numerator
of degree $13$. So $E_m$ satisfies for $m\ge14$ the linear recurrence of order
$12$ with characteristic polynomial $x^{12}L(1/x)^3N(1/x)^3$. Its roots are
$\rho$, $\sigma$, and $\pm i\sqrt2$, each of multiplicity $3$. The right side
of~\eqref{eq:Eclosed} is a combination of $m^j\rho^m$, $m^j\sigma^m$, and
$m^j(\pm i\sqrt2)^m$ with $j\le2$, so it satisfies the same recurrence for all
$m$. The two sides agree for $2\le m\le13$, which we checked in exact
arithmetic. Hence they agree for all $m\ge2$.

We have $|\sigma|<1$ and $|(-2)^{\lfloor m/2\rfloor}\kappa(m)|=O(m^22^{m/2})$,
and $\sqrt2<\rho$. This gives the asymptotic form of $E_m$. The expansion of
$a_3(m)$ follows, since
$4(m-1)(3m+2)3^{m-5}=\frac4{81}m^23^m-\frac4{243}(m+2)3^m$. Finally $E_m>0$
because $Xa^{m-1}$ is an optimum by Corollary~\ref{cor:threevalue}, and
$E_2=10$.
\end{proof}

\subsection{Optima up to symmetry}

Let $o_3(m)$ be the number of optima of the $3\times m$ grid up to the row
reflection and the column reversal. Its first values for $m=3,\dots,8$ are
$15$, $58$, $227$, $864$, $3190$, and $11502$ (\seqnum{A401226}).

\begin{theorem}[three-row orbits]\label{thm:threeorbits}
Let $M=1-2z^2-2z^4$. Then
\begin{align*}
4\sum_{m\ge3}o_3(m)z^m={}&\sum_{m\ge3}a_3(m)z^m+z^3+\frac{4z^3(1+z^4+z^6)}{M\,(1+2z^4)}\\
&+\frac{2z^3(1-z)(1+5z+3z^2-3z^3)}{(1-3z^2)^2}.
\end{align*}
For even $m\ge4$ we have $o_3(m)=\frac14\bigl(a_3(m)+2m\,3^{m/2-2}\bigr)$.
\end{theorem}

\begin{proof}
By Burnside's lemma, $4o_3(m)$ is $a_3(m)$ plus the numbers of optima fixed by
the row reflection, the column reversal, and the half turn.

\emph{Row reflection.} A fixed word uses only $b$, $B$, and $X$. A candidate of
this kind is $b^iXb^j$. For $i\ge2$ the window $bbX$ fails, and for $j\ge2$ the
window $Xbb$ fails (Table~\ref{tab:windows}). So the only fixed optimum is
$bXb$, which gives the term $z^3$.

\emph{Column reversal.} A fixed word is a palindrome. A candidate has
different doubles, so a fixed candidate has the form $UXU^{\mathrm{rev}}$ and
$m$ is odd. The window centered on $X$ requires the last letter of $U$ to be
$b$. As in Proposition~\ref{prop:GX}, the allowed strings $U$ are counted by
$v$. This gives $z\,v(z^2)$.

\emph{Half turn.} It reverses the word and exchanges $a$ with $c$ and $A$ with
$C$. A fixed word of type~I is $UXU'$, where $U'$ is the image of $U$. The
neighbors of $X$ are then $ac$, $ca$, or $bb$, which are the allowed pairs. So
these words give $z\,(2u+v)(z^2)$. A fixed word of type~II has $A$ and $C$ in
mirror columns $j$ and $m+1-j$. Put $m=2k$ or $m=2k+1$, where in the second
case the middle letter is $b$. Take the class $AC$, so $j\le k$. The first $k$
letters determine the word. Only the windows centered on $A$ and $C$ can fail,
and these $2$ windows are mirror images. If $j=1$, the window on $A$ does not
exist or does not restrict the first $k$ letters, and there are $3^{k-1}$
words. If $j=k\ge2$, the right neighbor of $A$ is $C$ or the middle $b$, so the
left neighbor is $a$ and there are $3^{k-2}$ words. If $1<j<k$, there are
$6$ allowed pairs of neighbors and $2\cdot3^{k-2}$ words. The sum over $j$ is
$2k\,3^{k-2}$ for $k\ge2$. The class $CA$ has the same count. So type~II
gives $2$ fixed optima for $m=3$ and $4k\,3^{k-2}$ for $m\ge4$, with series
\[
2z^3+\tfrac49(1+z)\Bigl(\frac{3z^2}{(1-3z^2)^2}-3z^2\Bigr)
=\frac{2z^3(1-z)(1+5z+3z^2-3z^3)}{(1-3z^2)^2}.
\]

By~\eqref{eq:hsol} and~\eqref{eq:runs} the $2$ contributions of the words
$UXU^{\mathrm{rev}}$ and $UXU'$ add up to
$2z\,(u+v)(z^2)=4z^3(1+z^4+z^6)/\Psi(z^2)$, and $\Psi(z^2)=M(1+2z^4)$. For even
$m\ge4$ only the type~II words are fixed, and $4k\,3^{k-2}=2m\,3^{m/2-2}$.
\end{proof}
